% =====================================================================
% Paper 55 (standalone): Periods of GeoVac
%
% Cyclotomic mixed-Tate classification of the master Mellin engine on
% the discrete Camporesi-Higuchi S^3 spectral triple. Companion synthesis
% paper to Papers 18, 28, 32, 38, 50.
%
% Audience: math-phys + NCG-with-motivic-leanings (Marcolli/Fathizadeh
% lineage, Brown/Glanois/Deligne periods-of-MZV program,
% Connes-Marcolli renormalization-motives).
%
% Status: first draft, June 2026, post-Sprints M1/M2/M3 (2026-06-03)
% which closed the master Mellin engine trinity at the Paper 32 §VIII
% corollary level.
%
% Built from internal sprint memos:
%   debug/sprint_mixed_tate_test_memo.md (M2)
%   debug/sprint_m3_cyclotomic_mixed_tate_memo.md (M3)
%   debug/sprint_m1_pure_tate_memo.md (M1)
% and the closure paragraphs in Paper 18 §III.7 and Paper 32 §VIII
% corollaries cor:m1_pure_tate, cor:m2_mixed_tate,
% cor:m3_cyclotomic_mixed_tate.
% =====================================================================
\documentclass[11pt,a4paper]{article}

\usepackage[T1]{fontenc}
\usepackage[utf8]{inputenc}
\usepackage{lmodern}
\usepackage{amsmath, amssymb, amsthm, mathrsfs, mathtools}
\usepackage{geometry}
\geometry{margin=1in}
\usepackage{hyperref}
\hypersetup{colorlinks=true, linkcolor=blue!50!black, citecolor=blue!50!black, urlcolor=blue!50!black}
\usepackage[expansion=false]{microtype}
\usepackage{graphicx}
\usepackage{booktabs}
\usepackage[numbers,sort&compress]{natbib}
\usepackage{xcolor}

\theoremstyle{plain}
\newtheorem{theorem}{Theorem}[section]
\newtheorem{lemma}[theorem]{Lemma}
\newtheorem{proposition}[theorem]{Proposition}
\newtheorem{corollary}[theorem]{Corollary}
\theoremstyle{definition}
\newtheorem{definition}[theorem]{Definition}
\newtheorem{remark}[theorem]{Remark}
\newtheorem{example}[theorem]{Example}

\newcommand{\nmax}{n_{\max}}
\newcommand{\sthree}{S^{3}}
\newcommand{\sfive}{S^{5}}
\newcommand{\R}{\mathbb{R}}
\newcommand{\C}{\mathbb{C}}
\newcommand{\N}{\mathbb{N}}
\newcommand{\Z}{\mathbb{Z}}
\newcommand{\Q}{\mathbb{Q}}
\newcommand{\Vol}{\mathrm{Vol}}
\newcommand{\Tr}{\mathrm{Tr}}
\newcommand{\MT}{\mathcal{MT}}
\newcommand{\Hcal}{\mathcal{H}}
\newcommand{\Mcal}{\mathcal{M}}
\newcommand{\Bcal}{\mathcal{B}}
\newcommand{\Zcal}{\mathcal{Z}}
\newcommand{\DCH}{D_{\mathrm{CH}}}

\title{Periods of GeoVac:\ a cyclotomic mixed-Tate classification of
the master Mellin engine}
\author{J.~Loutey\thanks{Independent researcher.
\texttt{jloutey@gmail.com}.  This work was developed in an
AI-augmented agentic workflow with Anthropic's Claude.
Implementation, internal sprint memos, and bibliographic collation
were performed collaboratively; all theorems, design choices, and
editorial decisions are the author's responsibility.}}
\date{June 2026 (preprint)}

\begin{document}
\maketitle

\begin{abstract}
The GeoVac framework is a discrete spectral-graph approach to
quantum mechanics built around the Fock conformal projection
$\R^3 \hookrightarrow \sthree$.  Its computational observables
empirically produce a constrained class of transcendental numbers:\
$\pi$ and integer powers thereof, classical odd zeta values
$\zeta(2k+1)$, Catalan $G = \beta(2)$ and the Dirichlet beta values
$\beta(s) = L(s, \chi_{-4})$, and a depth-$k$ tower of irreducible
constants at higher loop order (Paper~28 of the GeoVac corpus).  We
prove that every such transcendental sits in cyclotomic mixed-Tate
periods at level $\le 4$ over $\mathcal{O}_4[1/4] = \Z[i, 1/2]$, with
a complete period-theoretic stratification indexed by the
operator-order slot $k \in \{0, 1, 2\}$ of the master Mellin engine
$\Mcal[\Tr(D^k e^{-tD^2})]$ on the discrete Camporesi--Higuchi
$\sthree$ spectral triple:
\begin{itemize}
  \item $k = 0$ (Hopf-base measure $\Vol(S^2)/4 = \pi$):\
    $M_1 \subset \Q[\pi, \pi^{-1}]$, the localised pure-Tate algebra
    at $\pi$, depth $0$.
  \item $k = 2$ (Seeley--DeWitt coefficients):\
    $M_2 \subset \bigoplus_k \pi^{2k}\cdot\Q$, the pure-Tate
    even-weight sub-ring, strictly smaller than the generic
    Fathizadeh--Marcolli~\cite{fathizadeh_marcolli2016} mixed-Tate
    classification of Robertson--Walker spectral actions:\ on the
    static $\sthree$ sub-case, the Dirac SD expansion is two-term
    exact and the scalar Laplacian SD coefficients form a closed
    geometric series.
  \item $k = 1$ (vertex parity / Hurwitz at quarter-integer shifts /
    Dirichlet $L$):\
    $M_3 \subset \MT(\Z[i, 1/2])$ at level $\le 4$ via
    Deligne~\cite{deligne2010} and the explicit basis of
    Glanois~\cite{glanois2015}, at depth equal to loop order.  The
    vertex-parity-as-$\chi_{-4}$ correspondence is the
    Deligne--Glanois Galois descent from level $4$ to level $2$
    realised on the framework's natural QED observables.
\end{itemize}
The three sub-mechanisms are not set-theoretically disjoint
(joint engagement is the typical case in physical observables);
however, the master Mellin index $k$ uniquely identifies the
period-ring contribution of each spectral object.  We catalogue
witnesses from across the corpus (Papers 2, 25, 28, 38, 40, 43, 50,
51) and indicate open questions:\ specific motivic identification of
the irreducible depth-$k$ tower constants (PSLQ-irreducible at 200
dps starting from $S_{\min}$ at $k=2$), inhomogeneous extension to
non-static Robertson--Walker substrates, cross-manifold extension to
the Bargmann--Segal Hardy lattice on $\sfive$, and inner-factor
Yukawa coupling content.
\end{abstract}

\tableofcontents

% =====================================================================
\section{Introduction}
\label{sec:intro}
% =====================================================================

\subsection{Background:\ GeoVac and the transcendental-to-periods question}

The GeoVac framework~\cite{loutey_paper7,loutey_paper18} treats the
$\sthree$ conformal projection of the hydrogen Coulomb problem
(Fock 1935~\cite{fock1935}) as a discrete spectral-graph object:\
nodes labeled by quantum numbers $(n, l, m, \ldots)$, edges by
spectral selection rules, eigenvalues of the graph Laplacian giving
the Coulomb spectrum exactly via the universal scaling $\kappa = -1/16$.
The framework has been extended into a multi-paper corpus covering
quantum chemistry, qubit encodings, gauge structure, Lorentzian
extension, and gravity.  A consistent feature throughout:\ the
\emph{algebraic skeleton} of the framework (the bare graph and its
combinatorial selection rules) produces matrix elements in $\Q$, while
\emph{transcendentals} enter at projection steps that promote the
discrete substrate to continuum observables.  Paper~18 of the
corpus~\cite{loutey_paper18} catalogues these as ``exchange
constants'' organised by six tiers (intrinsic, calibration,
embedding, algebraic-implicit, composition, inner-factor input data).

The empirical content of the catalogue is short:\ across the entire
GeoVac corpus, only the following transcendentals have ever appeared
in spectral data:\
\begin{equation}
\label{eq:empirical_transcendentals}
\pi^{n}\,(n \in \Z), \quad
\zeta(2k+1)\,(k \ge 1), \quad
G = \beta(2), \quad
\beta(4),\;\beta(s),\;\ldots,
\quad\text{plus a depth-}k\text{ tower}\quad
S_{\min}, S^{(3)}, S^{(4)}, \ldots
\end{equation}
The first three classes have classical period-theoretic
identifications (Kontsevich--Zagier 2001~\cite{kontsevich_zagier2001};
Brown 2012~\cite{brown2012}; Deligne 2010~\cite{deligne2010};
Glanois 2015~\cite{glanois2015}).  The depth-$k$ tower constants are
PSLQ-irreducible at 200 dps against 100-element bases of all known
depth-$\le 2$ transcendentals (Paper~28 \S~$S_{\min}$,
\cite{loutey_paper28}); they are genuinely new transcendentals of the
framework.

The natural question, opened by the empirical content of
Eq.~\eqref{eq:empirical_transcendentals}:\ do all of GeoVac's
transcendentals sit in a single, named period-theoretic class?  If so,
what is the structural mechanism that selects this class?

\subsection{Result:\ master Mellin engine partition is a complete
period-ring stratification}

We show that the answer is affirmative.  Every transcendental in
Eq.~\eqref{eq:empirical_transcendentals} lies in cyclotomic mixed-Tate
periods at level $\le 4$ over $\mathcal{O}_4[1/4] = \Z[i, 1/2]$, with
a complete stratification indexed by the operator-order slot
$k \in \{0, 1, 2\}$ of the \emph{master Mellin engine}
\begin{equation}
\label{eq:master_mellin_engine}
\Mcal\!\left[\,\Tr\bigl(D^k\, e^{-tD^2}\bigr)\,\right](s)
\end{equation}
on the discrete Camporesi--Higuchi $\sthree$ spectral
triple~\cite{camporesi_higuchi1996}.  The structure is:
\begin{center}
\begin{tabular}{@{}clll@{}}
\toprule
$k$ & Sub-mechanism & Period-ring home & Source \\
\midrule
$0$ & M1 (Hopf-base measure) & $\Q[\pi, \pi^{-1}]$ & this paper \S\ref{sec:m1} \\
$2$ & M2 (Seeley--DeWitt) & $\bigoplus_k \pi^{2k}\cdot\Q$ pure Tate & \cite{fathizadeh_marcolli2016} + this paper \S\ref{sec:m2} \\
$1$ & M3 (vertex parity / Hurwitz / Dirichlet $L$) & $\MT(\Z[i, 1/2])$ at level $\le 4$ & \cite{deligne2010, glanois2015} + this paper \S\ref{sec:m3} \\
\bottomrule
\end{tabular}
\end{center}

The M3 row contains the substantive new identification:\ the
vertex-parity selection rule that produces Catalan $G$ and
$\beta(s)$ values in the framework's two-loop QED sector (Paper~28
Theorem~3, \cite{loutey_paper28}) is precisely the Deligne--Glanois
Galois descent from level $4$ to level $2$ on the framework's
natural Dirichlet series $D(s) = D_{\mathrm{even}}(s) +
D_{\mathrm{odd}}(s)$, with Hurwitz at quarter-integer shifts
$\zeta(s, 1/4), \zeta(s, 3/4)$ surviving in
$D_{\mathrm{even}} - D_{\mathrm{odd}}$ but cancelling in the sum.

\subsection{Structure}

Section~\ref{sec:master_engine} sets up the master Mellin engine on the
Camporesi--Higuchi spectrum and states the case-exhaustion theorem of
Paper~32 \S~VIII~\cite{loutey_paper32}.  Sections~\ref{sec:m1},
\ref{sec:m2}, and~\ref{sec:m3} prove the three sub-mechanism
classifications, each followed by a witness catalogue from the corpus.
Section~\ref{sec:joint} treats joint engagement and physical observables
that engage multiple sub-mechanisms simultaneously (the canonical
example being the K-formula $K = \pi(B + F - \Delta)$ of
Paper~2~\cite{loutey_paper2}).  Section~\ref{sec:open} catalogues
open questions:\ specific motivic identification of the depth-$k$
tower, inhomogeneous extension, cross-manifold extension to the
$\sfive$ Bargmann--Segal Hardy lattice~\cite{loutey_paper24}, and
inner-factor Yukawa content~\cite{loutey_paper32}.
Section~\ref{sec:conclusion} concludes.

% =====================================================================
\section{The master Mellin engine on Camporesi--Higuchi $\sthree$}
\label{sec:master_engine}
% =====================================================================

\subsection{The Camporesi--Higuchi spectrum}

Let $\DCH$ denote the canonical Dirac operator on the round $\sthree$
of unit radius, with the Camporesi--Higuchi spectral
realisation~\cite{camporesi_higuchi1996}:
\begin{equation}
\label{eq:ch_spectrum}
|\lambda_n^{\mathrm{CH}}| = n + \tfrac{3}{2}, \qquad
g_n = 2(n+1)(n+2), \qquad n \in \N_{\ge 0},
\end{equation}
where the half-integer spectrum arises from the spinor parallel
transport on the parallelisable group manifold $\sthree = \mathrm{SU}(2)$.
The squared Dirac spectrum is the integer-coefficient quadratic
$\lambda_n^2 = (n + 3/2)^2$ with the same degeneracy.  The associated
spectral zeta function is
\begin{equation}
\label{eq:dirac_dirichlet}
D(s) := \sum_{n \ge 0} \frac{g_n}{|\lambda_n^{\mathrm{CH}}|^s}
      = 2\,\zeta(s - 2, 3/2) - \tfrac{1}{2}\,\zeta(s, 3/2),
\end{equation}
where $\zeta(s, a)$ is the Hurwitz zeta function, and the squared
spectral zeta is
\begin{equation}
\label{eq:dirac_squared_dirichlet}
\zeta_{D^2}(s) := \sum_{n \ge 0} \frac{g_n}{(\lambda_n^{\mathrm{CH}})^{2s}}
                = D(2s).
\end{equation}

\subsection{The master Mellin engine}

The Mellin transform of the trace
$\Tr\bigl(D^k\, e^{-tD^2}\bigr)$ extracts the spectral zeta value
$D(2s - k)$ (modulo $\Gamma$-factor normalisations):
\begin{equation}
\label{eq:mellin_extract}
\Mcal\!\left[\Tr\bigl(D^k\, e^{-tD^2}\bigr)\right](s)
= \frac{1}{\Gamma(s)} \int_0^\infty t^{s-1} \,
   \sum_{n \ge 0} g_n (\lambda_n)^k\, e^{-t \lambda_n^2}\,dt
= \frac{\Gamma(s + k/2)}{\Gamma(s)} \cdot D(2s + k).
\end{equation}
At integer $s$, the three cases $k \in \{0, 1, 2\}$ produce
structurally distinct outputs:
\begin{itemize}
  \item $k = 0$:\ trivial heat kernel collapses to volume integral over
    the substrate; the leading $t \to 0^+$ coefficient is
    $\Vol(\sthree)/(4\pi)^{3/2}$ via the Weyl asymptotic.  The
    Hopf-bundle decomposition $\sthree \to S^2$ (Paper~25,
    \cite{loutey_paper25}) gives the base measure factor
    $\Vol(S^2)/4 = \pi$ in the explicit identifications below;\ we
    designate this contribution as mechanism M1.
  \item $k = 2$:\ standard Connes--Chamseddine spectral action
    $\Tr\,f(D^2/\Lambda^2)$ produces Seeley--DeWitt coefficients
    $\{a_k(D^2)\}_{k \ge 0}$ via the asymptotic expansion of the heat
    trace as $t \to 0^+$.  We designate this contribution as M2.
  \item $k = 1$:\ the $\eta$-invariant series
    $\eta(s) := \sum_n \mathrm{sgn}(\lambda_n) g_n |\lambda_n|^{-s}$
    arises by Mellin transform of $\Tr(D e^{-tD^2})$, with the
    vertex-parity decomposition
    $D = D_{\mathrm{even}} + D_{\mathrm{odd}}$ on the parity of
    $n_{\mathrm{Fock}}$ exposing Hurwitz at quarter-integer shifts.
    We designate this contribution as M3.
\end{itemize}

\subsection{The case-exhaustion theorem}

Paper~32 \S~VIII Theorem~\ref{thm:case_exhaustion}
(\cite{loutey_paper32}) establishes:
\begin{theorem}[Case-exhaustion of $\pi$-sources on the discrete
$\sthree$ spectral triple]
\label{thm:case_exhaustion}
Every appearance of $\pi$ in any closed-form GeoVac spectral
observable arises from one (or more) of the three named sub-mechanisms
$M_1, M_2, M_3$:
\begin{itemize}
  \item $M_1$:\ Hopf-base measure factor $\Vol(S^2)/4 = \pi$ on the
    Hopf bundle $\sthree \to S^2$;
  \item $M_2$:\ Seeley--DeWitt heat-kernel coefficient
    $a_k(D^2) \in \pi\cdot\Q$ in the asymptotic expansion of
    $\Tr\,e^{-tD^2}$;
  \item $M_3$:\ vertex-restricted Dirichlet-$L$ content via Hurwitz at
    quarter-integer shifts in the parity decomposition of $D(s)$.
\end{itemize}
A falsification sweep against three candidate fourth mechanisms
(continuum gravity, anomaly classes, principal bundles) returned the
sharper finding that $M_1, M_2, M_3$ are three sub-cases of the single
master Mellin engine $\Mcal[\Tr(D^k e^{-tD^2})]$ parameterised by
operator order $k \in \{0, 1, 2\}$.
\end{theorem}

The theorem partitions GeoVac's $\pi$-content by Mellin slot;\ it does
not by itself classify the \emph{period-theoretic ring} in which each
sub-mechanism's output lives.  That classification --- which is the
content of the present paper --- requires separate motivic input for
each slot, drawn from the published period theory of
Kontsevich--Zagier~\cite{kontsevich_zagier2001},
Fathizadeh--Marcolli~\cite{fathizadeh_marcolli2016},
Deligne~\cite{deligne2010}, and Glanois~\cite{glanois2015}.

% =====================================================================
\section{M1:\ the Hopf-base measure sub-mechanism}
\label{sec:m1}
% =====================================================================

\subsection{Period-ring statement}

\begin{theorem}[M1 period-ring]
\label{thm:m1_pure_tate}
The M1 sub-mechanism on the discrete Camporesi--Higuchi $\sthree$
spectral triple produces period values in the localised pure-Tate
sub-ring
\begin{equation}
\label{eq:m1_ring}
M_1 \;\subset\; \Q[\pi, \pi^{-1}]
\end{equation}
of mixed-Tate periods over $\Q$ (Kontsevich--Zagier
2001~\cite{kontsevich_zagier2001}), at depth $0$ and Tate weight
$\in \Z$.
\end{theorem}

\begin{proof}
By construction, the M1 mechanism is the Hopf-base measure factor
$\Vol(S^2)/4 = \pi$ on the Hopf bundle $\sthree \to S^2$
(Paper~25~\S~II.1,~\cite{loutey_paper25}).  As a period in the
Kontsevich--Zagier framework, $\pi$ has the canonical representation
$\pi = \tfrac{1}{2i}\oint_{|z|=1} dz/z$, a Tate weight-$1$ period of
the mixed Tate motive $\Q(1)$.  Every M1 closed form across the
corpus is a $\Q$-linear combination of integer powers of this period
(witnesses tabulated below);\ no M1 closed form has produced
transcendentals outside $\Q[\pi, \pi^{-1}]$.  The depth-$0$
property follows from the absence of iterated integration in the
Hopf-measure step:\ M1 acts as a single multiplicative factor, not as
a Mellin convolution.  The localisation at $\pi^{-1}$ accommodates the
inverse-period witnesses $4/\pi$ and $1/\pi^2$, which are not
periods in the strict Kontsevich--Zagier sense but sit naturally in
the field generated by $\pi$.
\end{proof}

\subsection{Witnesses}

\begin{center}
\begin{tabular}{@{}llcc@{}}
\toprule
Witness & Source & M1 factor & Tate weight \\
\midrule
$\Vol(S^2)/4 = \pi$ & Paper~25~\S~II.1~\cite{loutey_paper25} (Hopf base) & $\pi$ & $+1$ \\
$4/\pi$ (asymptotic propinquity rate) & Paper~38~\S~L2~\cite{loutey_paper38} & $\pi^{-1}$ & $-1$ \\
$4/\pi$ (universal across compact Lie groups) & Paper~40 Thm 1~\cite{loutey_paper40} & $\pi^{-1}$ & $-1$ \\
$1/\pi^2$ (Pythagorean orthogonality prefactor) & Paper~43~\S~10.2~\cite{loutey_paper43} & $\pi^{-2}$ & $-2$ \\
$1/\pi^2$ (F-theorem cross-product) & Paper~50 Thm 3.4~\cite{loutey_paper50} & $\pi^{-2}$ & $-2$ \\
\bottomrule
\end{tabular}
\end{center}

Every entry is $q \cdot \pi^n$ for $q \in \Q$ and $n \in \Z$ with
$|n| \le 2$ on the empirical witnesses logged so far.  The K-formula
$K = \pi(B + F - \Delta)$ of Paper~2~\cite{loutey_paper2} contains the
canonical M1 factor (the leading $\pi$);\ the analysis of joint
engagement in \S\ref{sec:joint} below decomposes the full M1+M2+M3
content of the K-formula.

\subsection{Convention-freeness}

Unlike M2 (which depends on the volume-normalisation convention for
the spectral action --- see \S\ref{sec:m2} below), the M1 ring is
convention-free.  The reason:\ the Hopf-base measure is itself a
volume integral, not a heat-kernel asymptotic coefficient at $t \to 0^+$.
There is no choice of dimensional rescaling to absorb or expose; the
factor $\pi$ appears directly.  This convention-freeness is the
structural counterpart of M1's depth-$0$ status:\ depth $\ge 1$
mechanisms admit non-trivial rescaling choices via the heat-kernel
asymptotic, while depth-$0$ mechanisms do not.

% =====================================================================
\section{M2:\ the Seeley--DeWitt sub-mechanism}
\label{sec:m2}
% =====================================================================

\subsection{Period-ring statement}

\begin{theorem}[M2 period-ring]
\label{thm:m2_mixed_tate}
The M2 sub-mechanism on the discrete Camporesi--Higuchi $\sthree$
spectral triple produces period values in the strictly smaller
pure-Tate even-weight sub-ring
\begin{equation}
\label{eq:m2_ring}
M_2 \;\subset\; \bigoplus_{k \ge 0} \pi^{2k}\cdot\Q
\;\subset\; \MT(\Q)
\end{equation}
of mixed-Tate periods over $\Q$, at depth $0$.  The classification
inherits the Fathizadeh--Marcolli~\cite{fathizadeh_marcolli2016}
mixed-Tate result for Connes--Chamseddine spectral actions on
Robertson--Walker spacetimes $\R \times \sthree$, restricted to the
static sub-case $a(t) \equiv 1$, and refines it in two structural
directions:\ (i) the Dirac SD expansion is two-term exact ($a_k = 0$
for $k \ge 2$), while the F--M generic R--W expansion is infinite;\
(ii) the scalar Laplacian SD coefficients form the closed
geometric series $a_k^\Delta = 2\pi^2/k!$, while the F--M generic
R--W scalar expansion has no such closed form.
\end{theorem}

\begin{proof}[Proof sketch]
Direct inheritance is by application of
Fathizadeh--Marcolli~\cite{fathizadeh_marcolli2016} Theorem on the
Rosenfeld integral representation of Seeley--DeWitt coefficients on
$\R \times \sthree$, specialised to constant scaling factor.  The
pure-Tate refinement on the static sub-case has two parts:
\begin{itemize}
  \item \textit{Dirac sector.}  The Camporesi--Higuchi spectrum
    $|\lambda_n| = n + 3/2$ admits the Jacobi $\vartheta_2$ modular
    transformation
    $\Tr e^{-tD^2} = (\sqrt\pi/2) t^{-3/2} - (\sqrt\pi/4) t^{-1/2}
                    + O(e^{-\pi^2/t})$, with the exponentially small
    remainder vanishing in the asymptotic series.  All SD coefficients
    $a_k^{D^2}$ for $k \ge 2$ are exactly zero;\ the volume-normalised
    coefficients are $a_0 = 4\pi^2$, $a_1 = -2\pi^2$.
  \item \textit{Scalar sector.}  The integer spectrum $n(n+2)$ admits
    the Jacobi $\vartheta_3$ modular transformation giving the closed
    form $\Tr e^{-t\Delta} = (\sqrt\pi/4) \cdot e^t/t^{3/2} +
    O(e^{-\pi^2/t})$, with $a_k^\Delta = 2\pi^2/k!$.
\end{itemize}
Both sectors are explicit in Paper~51~\cite{loutey_paper51}.  The
empirical pure-Tate restriction (no $\zeta(3)$, no MZVs) is consistent
with this closed-form structure:\ the F--M generic R--W classification
allows MZV content when scaling-factor derivatives are non-trivial;\
on the static sub-case those derivatives vanish, eliminating the MZV
content.
\end{proof}

\subsection{Witnesses}

The empirical pure-Tate witness panel from the GeoVac corpus:

\begin{center}
\begin{tabular}{@{}llcc@{}}
\toprule
Witness & Source & Closed form & Tate weight \\
\midrule
$a_0^{D^2}$ (Dirac SD, vol-norm) & Paper~51 Cor.~2.1~\cite{loutey_paper51} & $4\pi^2$ & $+2$ \\
$a_1^{D^2}$ (Dirac SD, vol-norm) & Paper~51 Cor.~2.1 & $-2\pi^2$ & $+2$ \\
$a_k^{D^2}$, $k \ge 2$ & Paper~51 Cor.~2.1 & $0$ & --- \\
$a_k^\Delta$ (scalar SD, vol-norm) & Paper~51 Thm 3.1 & $2\pi^2/k!$ & $+2$ \\
$\zeta_{D^2}(2)$ & Paper~28 Thm 1~\cite{loutey_paper28} & $\pi^2 - \pi^4/12$ & $+2, +4$ \\
$\zeta_{D^2}(3)$ & Paper~28 Thm 1 & $\pi^4/3 - \pi^6/30$ & $+4, +6$ \\
$\zeta_{D^2}(4)$ & Paper~28 Thm 1 & $2\pi^6/15 - 17\pi^8/1260$ & $+6, +8$ \\
$\zeta_{D^2}(5)$ & Paper~28 Thm 1 & $\pi^8(306 - 31\pi^2)/5670$ & $+8, +10$ \\
\bottomrule
\end{tabular}
\end{center}

All witnesses are in $\bigoplus_k \pi^{2k}\cdot\Q$;\ no $\sqrt\pi$,
no $\zeta(2k+1)$, no MZV content at the SD or spectral-zeta-at-integer-$s$
level.

\subsection{Volume-normalisation convention}

A structural caveat at the raw heat-trace level:\ the un-normalised
coefficients $\tilde a_k = a_k / (4\pi)^{3/2}$ contain $\sqrt\pi$
explicitly (from the spinor fiber dimension factor that always appears
in $d$-odd heat-kernel expansions on spin manifolds).  This $\sqrt\pi$
cancels exactly against the $(4\pi)^{d/2}$ in the standard
volume-normalisation convention, recovering the pure-Tate ring
$\bigoplus_k \pi^{2k}\cdot\Q$.  At the spectral-zeta-at-integer-$s$
level (which is convention-free), the values lie in
$\Q[\pi^2]$;\ the $\sqrt\pi$ slot is populated only at half-integer
$s$, where the values match the leading heat-trace poles.  This
reconciles Paper~28 Theorem~1 (T9, $\sqrt\pi\cdot\Q \oplus \pi^2\cdot\Q$
two-term ring) with the pure-Tate classification at integer $s$.

% =====================================================================
\section{M3:\ vertex parity, Hurwitz at quarter-integer shifts, and
the Deligne--Glanois Galois descent}
\label{sec:m3}
% =====================================================================

[\textbf{Drafting in progress.}  This section will state and prove
the M3 period-ring theorem, parallel to Theorems~\ref{thm:m1_pure_tate}
and~\ref{thm:m2_mixed_tate}, with full witness panel from Paper~28
Theorems 2--3 ($D(s)$ parity discriminant; $\chi_{-4}$ identity),
Paper~28 §K-Sommerfeld separation (Sommerfeld $D_p$ with harmonic-number
weights), Paper~28 §$S_{\min}$ (two-loop irreducible at 200 dps), and
Paper~28 §depth-$k$ tower.  The substantive new content of \S\ref{sec:m3}
is the identification of the vertex-parity-as-$\chi_{-4}$ correspondence
as the literal Deligne--Glanois Galois descent at $N = 4$ on the
framework's natural QED observables.  Source memo:
\texttt{debug/sprint\_m3\_cyclotomic\_mixed\_tate\_memo.md}.]

% =====================================================================
\section{Joint engagement and physical observables}
\label{sec:joint}
% =====================================================================

[\textbf{Outline.}  Sub-mechanisms M1, M2, M3 are role-disjoint at the
Mellin slot index but not set-theoretically disjoint;\ joint engagement
is the typical case in physical observables.  Canonical examples:
\begin{itemize}
  \item The K-formula $K = \pi(B + F - \Delta)$ of
    Paper~2~\cite{loutey_paper2}:\ explicit $\pi$ from M1, $F = \zeta(2)$
    from M2 inside the Fock Dirichlet evaluation, $\Delta = 1/40$
    rational from M3 weight-$0$ Dirac mode count.
  \item Paper~50 F-theorem coefficient
    $F_s = (\log 2)/8 - 3\zeta(3)/(16\pi^2)$:\ $1/\pi^2$ from M1,
    $\zeta(3)$ from M3 in cross-product.  The decomposition
    $F_D + 2F_s = (\log 2)/2$ isolates M2 (all $\zeta(3)$ cancels);
    the decomposition $F_D - 2F_s = 3\zeta(3)/(4\pi^2)$ isolates M3
    (all $\log 2$ cancels).
  \item Paper~38 L2 asymptotic propinquity rate constant
    $4/\pi = \Vol(S^2)/\pi^2$:\ M1 factor $\Vol(S^2)$ in the
    numerator, M1 inverse-period $\pi^{-2}$ in the denominator.
\end{itemize}
Section will formalise the joint-engagement combinatorics as a
$(k_1, k_2, k_3) \to k_1 + k_2 + k_3$ product of Tate weights
modulo the master Mellin index multiplicativity.]

% =====================================================================
\section{Open questions}
\label{sec:open}
% =====================================================================

\subsection{Specific motivic identification of the depth-$k$ tower}

The depth-$k$ loop tower (Paper~28 Proposition depth\_k,
\cite{loutey_paper28}) predicts a new irreducible transcendental at
each loop order:
\begin{itemize}
  \item $k = 1$:\ $\pi^{\mathrm{even}}$ only (Paper~28 Theorem~T9,
    proven; depth $1$ in $\MT(\Z)$).
  \item $k = 2$:\ $S_{\min} = 2.479936938\ldots$ PSLQ-irreducible at
    200 dps against a 100-element basis (Paper~28 \S$S_{\min}$);\
    classification predicts $S_{\min} \in \MT(\Z[1/2])$ at depth $2$.
  \item $k = 3$:\ $S^{(3)}$ independent of $S_{\min}$ (PSLQ with
    $S_{\min}$ in basis fails at fast-converging exponents).
  \item $k$ general:\ new irreducible at depth $k$.
\end{itemize}
The classification places each constant in a specific period ring
(by inheritance of the motivic content of the Hurwitz / Dirichlet $L$
building blocks) but does not identify the constant as a specific
element of the Deligne--Glanois basis $\Bcal^N$ at depth $k$.  Higher
PSLQ precision (300 dps) against an enriched depth-$k$ basis is the
natural mini-sprint at each $k$.

\subsection{Inhomogeneous extension}

The Fathizadeh--Marcolli result holds on $\R \times \sthree$ with
arbitrary scaling factor $a(t)$;\ the present paper restricts to
$a(t) \equiv 1$.  Whether the GeoVac discrete substrate admits a
time-dependent $a(t)$ generalisation with the same period-theoretic
content is open.

\subsection{Cross-manifold extension to $\sfive$ Bargmann--Segal}

Paper~24~\cite{loutey_paper24} establishes a parallel discrete
spectral-triple structure on the Hardy sector of $\sfive$, encoding
the 3D harmonic oscillator with a $\pi$-free skeleton.  Whether
the master Mellin engine partition extends to this cross-manifold
substrate, and whether the period-theoretic classification holds at
the same level $4$, is open.

\subsection{Inner-factor Yukawa content}

The almost-commutative extension
$\mathcal{A}_{\mathrm{GV}} \otimes (\C \oplus \mathbb{H})$ studied in
Paper~32 \S~VIII.C~\cite{loutey_paper32} (Sprint H1 closure) admits a
Higgs structure but does not autonomously select the Yukawa coupling
matrix.  The associated inner-factor input data (Yukawa Dirichlet
ring, Paper~18 \S~IV tier 6~\cite{loutey_paper18}) is classified
empirically as Dirichlet-$L$-class, but its period-theoretic placement
inside the M1/M2/M3 trinity (or as a genuinely new sub-mechanism) is
open.

% =====================================================================
\section{Conclusion}
\label{sec:conclusion}
% =====================================================================

The transcendental content of GeoVac spectral data lies in
cyclotomic mixed-Tate periods at level $\le 4$ over $\Z[i, 1/2]$,
indexed by the master Mellin engine slot $k \in \{0, 1, 2\}$.  The
classification places the GeoVac framework inside the published
periods/motives program (Kontsevich--Zagier, Brown, Deligne, Glanois,
Fathizadeh--Marcolli) in a precise way:\ each of GeoVac's
empirically-observed transcendentals has a named motivic home
inherited from the period-theoretic content of the corresponding
sub-mechanism.  The vertex-parity-as-$\chi_{-4}$ correspondence at the
core of GeoVac's two-loop QED sector (Paper~28 Theorem~3) is the
literal Deligne--Glanois Galois descent from level $4$ to level $2$
on the framework's natural Dirichlet series.

% =====================================================================
\begin{thebibliography}{99}

\bibitem{loutey_paper2}
J.~Loutey, ``Paper 2: $\alpha = K^{-1}$ as a spectral coincidence on
the Fock-projected $\sthree$,'' GeoVac corpus (2026).

\bibitem{loutey_paper7}
J.~Loutey, ``Paper 7: The dimensionless vacuum on $\sthree$,'' GeoVac
corpus (2026).

\bibitem{loutey_paper18}
J.~Loutey, ``Paper 18: Exchange constants and the master Mellin
engine,'' GeoVac corpus (2026).

\bibitem{loutey_paper24}
J.~Loutey, ``Paper 24: Bargmann--Segal lattice for the 3D harmonic
oscillator,'' GeoVac corpus (2026).

\bibitem{loutey_paper25}
J.~Loutey, ``Paper 25: GeoVac Hopf graphs as discrete lattice-gauge
structure,'' GeoVac corpus (2026).

\bibitem{loutey_paper28}
J.~Loutey, ``Paper 28: QED on $\sthree$,'' GeoVac corpus (2026).

\bibitem{loutey_paper32}
J.~Loutey, ``Paper 32: The GeoVac spectral triple,'' GeoVac corpus
(2026).

\bibitem{loutey_paper38}
J.~Loutey, ``Latr\'emoli\`ere propinquity convergence of truncated
Camporesi--Higuchi spectral triples on $\sthree$,'' arXiv preprint /
GeoVac corpus Paper~38 (2026).

\bibitem{loutey_paper40}
J.~Loutey, ``Universal $4/\pi$ rate constant in Latr\'emoli\`ere
propinquity convergence across compact connected Lie groups,'' GeoVac
corpus Paper~40 (2026).

\bibitem{loutey_paper43}
J.~Loutey, ``Lorentzian extension of the four-witness Wick-rotation
theorem at finite cutoff,'' GeoVac corpus Paper~43 (2026).

\bibitem{loutey_paper50}
J.~Loutey, ``Boundary CFT$_3$-on-$\sthree$ observables on the
truncated GeoVac spectral triple,'' GeoVac corpus Paper~50 (2026).

\bibitem{loutey_paper51}
J.~Loutey, ``Gravity arc on the discrete $\sthree$ spectral triple,''
GeoVac corpus Paper~51 (2026).

\bibitem{fock1935}
V.~Fock, ``Zur Theorie des Wasserstoffatoms,''
\textit{Z. Physik}~\textbf{98}, 145--154 (1935).

\bibitem{camporesi_higuchi1996}
R.~Camporesi and A.~Higuchi, ``On the eigenfunctions of the Dirac
operator on spheres and real hyperbolic spaces,''
\textit{J. Geom. Phys.}~\textbf{20}, 1--18 (1996).

\bibitem{kontsevich_zagier2001}
M.~Kontsevich and D.~Zagier, ``Periods,'' in \textit{Mathematics
Unlimited --- 2001 and Beyond}, B.~Engquist and W.~Schmid, eds.,
Springer (2001), pp.~771--808.

\bibitem{brown2012}
F.~Brown, ``Mixed Tate motives over $\Z$,''
\textit{Ann. of Math.}~(2)~\textbf{175}, 949--976 (2012).

\bibitem{deligne2010}
P.~Deligne, ``Le groupe fondamental unipotent motivique de
$\mathbb{G}_m - \mu_N$, pour $N = 2, 3, 4, 6$ ou $8$,''
\textit{Publ. Math. IHES}~\textbf{112}, 101--141 (2010);\ preprint
arXiv:math/0302267 (2005).

\bibitem{glanois2015}
C.~Glanois, ``Motivic unipotent fundamental groupoid of
$\mathbb{G}_m \setminus \mu_N$ for $N = 2, 3, 4, 6, 8$ and Galois
descents,'' arXiv:1411.4947 (v2, 2 Sep 2015);\ published version
\textit{J. Number Theory}~\textbf{182}, 36--90 (2018).

\bibitem{fathizadeh_marcolli2016}
F.~Fathizadeh and M.~Marcolli, ``Periods and motives in the spectral
action of Robertson--Walker spacetimes,''
\textit{Comm. Math. Phys.}~\textbf{356}, 641--671 (2017);\ preprint
arXiv:1611.01815 (2016).

\end{thebibliography}

\end{document}
